% !TEX root = lectures.tex
\section{Lecture 1}
When one develops the theory of machine learning, it's a good idea to follow the
following order:

\begin{center}
\begin{tikzpicture}
    \node (1) {Problem};
    \node (2) [right=of 1] {Theoretical Model};
    \node (3) [right=of 2] {Algorithmic Method};
    \node (4) [right=of 3] {Theoretical Guarantee};
    \draw[->] (1) edge (2);
    \draw[->] (2) edge (3);
    \draw[->] (3) edge (4);
\end{tikzpicture}
\end{center}
\begin{example}
    Here are some \textbf{problems} you may start with.
    \begin{enumerate}
        \item Prediction from expert advice
        \item Email spam classification
        \item Routing packets in a network
        \item Portfolio selection
        \item Robot control
        \item Sequence prediction
    \end{enumerate}
\end{example}

\subsection{Prediction From Expert Advice}
In this setup, we perhaps have a decision between two options $A$ and $B$. I
have $N$ experts giving advice on which option to pick. How do we choose the better option?
This kind of setup doesn't make much sense with a one-shot position; we have no information
on whether the experts give good advice.

How about the multi-shot setting? How can we hope to obtain the smallest number of mistakes?
In some sense we wish to identify those experts that are ``good.'' Well, if all the experts
are bad, then we can't really hope to do better than them. Instead,
a good model is to have the number of mistakes approach the best expert in hindsight.

Let's try a simple strategy to find this best expert. One simple strategy is to pick the expert that
has been the best so far. However, one can see quickly that this does not work. For instance, suppose
the predictions are as follows (T for true, F for false)

\begin{center}
\begin{tabular}{c | c}
    1 & 2 \\ \hline
    F & T \\
    T & F \\
    F & T \\
    T & F
\end{tabular}
\end{center}

Then it's clear that we will alternate between the experts and never get it right!
But the best expert in hindsight is right 50\% of the time. The following \textbf{weighted majority} 
algorithm does better.

\begin{algothm}[Littlestone/Warmuth '89]
    We will assign a weight to each experts $i$ at each time step $t$, call it $W_t(i)$.
    Initially, for all $i$, $W_1(i) = 1$. To make a decision at time $t$, we
    pick the decision $D$ that maximizes 
    \[ \sum_{\text{expert } i \text{ makes prediction } D \text{ at time } t} W_t(i) \]

    Every time step, we update the weight as follows. Take some parameter $\epsilon \in (0, 1)$.
    \[ W_{t + 1}(i) = \begin{cases}
        W_t(i) & \text{if } i \text{ was right} \\
        W_t(i) \cdot (1-\epsilon) & \text{ otherwise}
    \end{cases} \]
\end{algothm}

\begin{theorem}
    The number of mistakes made by the weighted majority,
    \[ M \le 2 (1 + \epsilon) R + \frac{\log N}{\epsilon} \]
    where $R$ is the number of mistakes made by the best expert.
    \begin{proof*}
        We will use a potential function argument to prove this.
        The potential function will be $\phi_t = \sum_{i = 1}^{N} W_t(i)$,
        the total sum of the weights. It's clear that $\phi_1 = N$ and $\phi_{t + 1} \le \phi_t$.
        In fact, if the algorithm errored, then at least half of the weight was on the wrong choice, so therefore
        we decrease at least half the weight. Thus
        \[ \phi_{t + 1} \le \frac{1}{2} \phi_t + \frac{1}{2} \phi_t (1 - \epsilon) \]
        Which means in fact a strict decrease of the weights, wherein $\phi_{t + 1} \le \phi_t\left(1 - \frac{\epsilon}{2}\right)$.
        Thus, $\phi_T \le \phi_1 \qty(1 - \frac{\epsilon}{2})^{M}$. Therefore, the highest weight
        is at most $N \qty(1 - \frac{\epsilon}{2})^{M}$.
        The highest weight is also $(1-\epsilon)^R$. We thus have
        \[(1-\epsilon)^R \le N \qty(1 - \frac{\epsilon}{2})^{M} \implies M \le  2 (1 + \epsilon) R + \frac{\log N}{\epsilon} \]
        by using $-x + x^2 \le \log(1 - x) \le - x$ for $0 < x < 1/2$.
    \end{proof*}
\end{theorem}

In fact, this algorithm is optimal for deterministic algorithms (as $\epsilon \to 0$).
The adversarial case is that the two experts always pick opposite things (fixed), and the
one your algorithm picks is always wrong. Then, one of the experts, by pigeonhole,
was right on at least half the days, but you were wrong every day.
It turns out we can get rid of this problem if we use a distribution.

\begin{algothm}[Randomized Weighted Majority]
    We make the decision $D$ with probability
    \[ \frac{\sum_{\text{expert } i \text{ makes prediction } D \text{ at time } t} W_t(i)}{\sum_{i} W_t(i)}  \] 
    The update and initial rules are the same.   
\end{algothm}

\begin{theorem}
    In the randomized weighted majority algorithm, $M$, follows
    \[ \E{M} \le (1 + \epsilon) R + \frac{\log N}{\epsilon} \]

    \begin{proof*}
    The proof is very similar. The only difference is now we can
    reason about $\E{\phi_t}$. Let $m_t$ be the indicator
    random variable on whether or not the algorithm made a mistake
    at time $t$. Let $m_t(i)$ be the indicator on whether expert $i$
    made a mistake (not random).
    Now, the new claim is $\phi_{t + 1} \le \phi_t(1 - \epsilon)^{\E{m_t}} \le e^{-\epsilon \E{m_t}}$.
    It's not hard to see that this will immediately yield the result we want.
    \begin{align*}
        \phi_{t + 1} &= \sum_i W_t(i) (1 - \epsilon m_t(i)) \\
        &= \phi_t (1 - \epsilon \sum_i p_t(i) m_t(i)) \\
        &= \phi_t (1 - \epsilon \E{m_t}) \le \phi_t e^{-\epsilon \E{m_t}}
    \end{align*}
\end{proof*}

\end{theorem}
